\section{Выбор средств реализации сервиса для программной поддержки генерации расписания учебных заведений}

\subsection{Обоснование выбора СУБД}

При выборе СУБД требуется учесть главные факторы:

\begin{itemize}
	\item размер данных: для генерации расписания не требуется огромный объём данных. Учебные планы, списки групп, списки преподавателей и т. д. занимают относительно небольшое место на диске;
	\item тип данных: данные для генерации расписаний являются текстовыми данными и легко могут храниться в любой СУБД или даже просто в текстовом файле, как JSON или CSV;
	\item требование к производительности: нагрузка на базу данных будет минимальная – генерировать расписание не требуется чаще, чем раз в семестр;
	\item простота использования: главной задачей ВКР является сам алгоритм генерации расписания, база данных является лишь второстепенной задачей, которая не должна отнимать много времени на развёртывание и заполнение.
\end{itemize}

Учитывая вышеперечисленные факторы, можно прийти к выводу, что для поставленной задачи требуется небольшая, легковесная база данных.

\subsubsection{Хранение данных в текстовых файлах}

Хранения данных в текстовых файлах, таких как JSON, TXT или CSV отлично подходит для поставленной задачи. Из плюсов можно выделить простоту использования, подходит для хранения сложных, вложенных структур данных, легко читается человеком и легко вносить изменения в структуру данных.

Однако есть и минусы. При каждом чтении будет производится операция парсинга текстовых данных, что будет делать операции чтения, изменения или добавления сравнительно медленными. Более того, такие, классические для СУБД вещи как поиск, индексация, транзакции отсутствуют, а сами данные будут храниться не в бинарном файле, а в текстовом, что будет неоправданно раздувать объём хранимых данных.

Текстовый формат данных отлично подходит для транспортировки небольших объёмов данных, таких как уже сгенерированное расписание, на клиент для дальнейшего отображения расписания учебных занятий. Однако текстовые файлы не идеально подходят для хранения и работы с данными, используемыми в процессе генерации расписания.

\subsubsection{SQLite}

SQLite – это самодостаточная и встроенная СУБД. Для развёртывания которой не требуется сервер или сложная система управления. Все данные хранятся в бинарном файле прямо в проекте с расширением “.sqlite” или “.db”.

Из минусов часто выделяют ограничение по объёму данных до 140Тб, однако это не будет проблемой для хранения данных для генерации расписаний учебных занятий. Также SQLite не подходит для сложных многопользовательских сценариев – отсутствует полноценная поддержка конкурентности.

\subsubsection{Вывод}

При небольшом анализе вариантов СУБД для хранения данных проекта выбор пал на SQLite. Это легковесная СУБД, обладающая всеми преимуществами классических СУБД как индексация и скорость обработки запросов, но и в то же время обладающая всеми преимуществами хранения данных в текстовых файлах – простота и гибкость. Более того, у автора данной выпускной квалификационной работы уже был опыт работы с SQLite, а соответственно не будет проблем с использованием, связанных с низким уровнем компетенции в инструменте.

\subsection{Выбор средств реализации сервера}

\subsection{Выбор средств реализации клиента}

\subsection{Вывод по четвёртой главе}

Был произведён сравнительный анализ подходящих для проекта СУБД и сделан обоснованный выбор на SQLite. 